\chapter{The Yang--Mills Gradient Flow Program}

In 2009--2010 Martin L\"uscher introduced a diffusion equation for gauge
fields and turned it into one of the most fruitful ideas of modern lattice
field theory. The flow smooths gauge fields over a radius $\sqrt{8t}$
while preserving gauge invariance, and the smoothed fields are
\emph{renormalized fields}: composite observables built from them are
finite without additional renormalization. The library contains the
original articles of this program; they anchor its two offshoots ---
precision scale setting and thermodynamics on one side, smeared operator
product expansions and smeared quasi-PDFs on the other.

% ----------------------------------------------------------------------
\section{Smooth Wilson Loops: The Idea of a Flow Before the Flow}
\papercard{R.\ Narayanan and H.\ Neuberger}
{Infinite N phase transitions in continuum Wilson loop operators}
{JHEP \textbf{03} (2006) 064}
{arXiv:hep-th/0601210 | DOI: 10.1088/1126-6708/2006/03/064}
\whyselected{An original article of this library and the conceptual
ancestor of gradient-flow smearing of Wilson lines: it is cited by the
Monahan--Orginos and Monahan quasi-PDF papers, by the Brambilla--Wang
off-lightcone paper, and by the higher-order flow papers ($\sim$6/99).}
\physics{A large Wilson loop is not a well-defined observable in the
continuum limit --- its perimeter diverges. Narayanan and Neuberger
studied a \emph{smoothed} loop obtained by diffusing the contour,
and showed numerically that the smoothed loop has a finite continuum
limit. In the 't~Hooft limit, enlarging the loop triggers a phase
transition where the eigenvalue density of the holonomy closes its gap;
if the transition falls in a solvable universality class (the
Gross--Witten / Durhuus--Olesen / Douglas--Kazakov lineage cited in the
paper), the string tension becomes analytically accessible.}
\impact{Beyond large-$N$ physics, this paper established that smoothing a
Wilson line is compatible with a finite continuum limit --- precisely the
fact on which smeared (pseudo-/quasi-) PDFs rest, since their nonlocal
operators contain open Wilson lines. The library's flow-based PDF papers
cite it as the origin of the idea.}
\cardend

% ----------------------------------------------------------------------
\section{Trivializing Maps: The Flow Enters Field Theory}
\papercard{M.\ L\"uscher}
{Trivializing maps, the Wilson flow and the HMC algorithm}
{Commun.\ Math.\ Phys.\ \textbf{293}, 899--911 (2010)}
{arXiv:0907.5491 | DOI: 10.1007/s00220-009-0995-5}
\whyselected{An original article of this library, and the first appearance
of the Yang--Mills gradient flow in the literature: the ``Wilson flow'' is
named here. Cited throughout the flow corpus of the library
($\sim$6/99).}
\physics{A trivializing map transforms the action to one whose
long-distance couplings vanish, ideally rendering the measure Gaussian.
L\"uscher shows that integrating the infinitesimal version leads to a
diffusion equation on field space,
\begin{equation}\label{eq:trivmap}
  \partial_s B_\mu(s,u) = E_\mu(B(s,u)), \qquad B_\mu(0)=U_\mu ,
\end{equation}
where $E_\mu$ is the negative $SU(3)$ gradient of the action --- the flow.
Exact maps would make Monte Carlo trivial; approximate ones can still be
used inside HMC to damp topological autocorrelations. The paper already
notes that short flow times probe the ultraviolet structure of the theory,
setting the agenda for everything that follows.}
\impact{This is where the community learned that flowing fields is both a
numerical tool and a renormalization concept. The AI-sampling papers of
Chapter~8 are, in a precise sense, machine-learning searches for
trivializing maps.}
\cardend

% ----------------------------------------------------------------------
\section{Properties of the Wilson Flow}
\papercard{M.\ L\"uscher}
{Properties and uses of the Wilson flow in lattice QCD}
{JHEP \textbf{08} (2010) 071}
{arXiv:1006.4518 | DOI: 10.1007/JHEP08(2010)071}
\whyselected{The single most-cited paper of the entire citation network
mined for this selection ($\sim$15/99): every gradient-flow article of the
library builds on it, and even LaMET and BMW papers quote it for scale
setting or QED-flow interpolation. An original article of the library.}
\physics{The Wilson flow evolves link fields by the gauge-covariant heat
equation
\begin{equation}\label{eq:flow}
  \partial_t B_\mu = D_\nu G_{\nu\mu} + \xi\,D_\mu\,\partial_\nu B_\nu,
  \qquad B_\mu(0)=A_\mu ,
\end{equation}
with $\xi=0$ (Wilson) or $1/5$ (symmetric). Combining one-loop perturbation
theory with numerical evidence, L\"uscher establishes three facts:
(i)~$B_\mu(t)$ at any fixed $t>0$ is a smooth field whose local gauge-invariant
functionals have finite continuum limits --- the flow is itself a
renormalization; (ii)~the observable $w(t)=t^2\langle E(t)\rangle$ with
$E=\frac12 G^a_{\mu\nu}G^a_{\mu\nu}$ rises steeply through a crossover at
$t\sim a^2$, so the scale $w_0$, defined by $\sqrt{8t}\,w(t)|_{w_0}=w_0$
($\approx0.4$~fm), provides a precision reference scale; (iii)~topological
sectors decouple dynamically once flowed beyond $\sqrt{8t}\sim a$.}
\impact{$w_0/w_a$ scale setting is standard practice across the library's
ensembles (BMW, CLQCD clover sets). More profoundly, the theorem
``flowed products are UV finite'' is the license for every application in
this chapter and the smeared-operator applications of Chapter~6.}
\cardend

% ----------------------------------------------------------------------
\section{Renormalizability of the Flowed Theory}
\papercard{M.\ L\"uscher and P.\ Weisz}
{Perturbative analysis of the gradient flow in non-abelian gauge theories}
{JHEP \textbf{02} (2011) 051}
{arXiv:1101.0963 | DOI: 10.1007/JHEP02(2011)051}
\whyselected{Second-highest hub of the flow corpus ($\sim$13/99): all
short-flow-time calculations in the library (energy--momentum tensor at
NNLO, quark bilinears through NNLO, dipole matching) cite it as their
axiomatic base. Original article of the library.}
\physics{The flow lives on the half-space $\mathbb{R}^4\times[0,\infty)$
with boundary condition $B(0)=A$; L\"uscher and Weisz formulate it as a
five-dimensional local field theory, derive its Feynman rules, and prove
by power counting plus BRST symmetry that correlation functions of flowed
fields at strictly positive flow time require \emph{no renormalization
beyond that of the underlying four-dimensional theory}. The proof
organizes divergences into counterterms that are themselves flows of
renormalized lower-dimension operators, with coefficients fixed recursively;
BRST Ward identities kill the dangerous classes. This is the all-orders
theorem on which every ``small-flow-time expansion'' rests:
\begin{equation}
  O_R(x) \;\xrightarrow{t\to0}\; c_O(t,\mu)\,[O]_R(x) + \mathcal{O}(t),
\end{equation}
with computable Wilson coefficients $c_O$ and finite, unambiguous limits.}
\impact{Every matched observable in this book's chapters 3 and 6 ---
running couplings, energy--momentum tensor, chiral condensate, smeared
twist-2 operators, smeared quasi-PDFs --- is an instance of this expansion
with coefficients computed within the five-dimensional formalism
introduced here.}
\cardend

% ----------------------------------------------------------------------
\section{Fermions and Chiral Symmetry}
\papercard{M.\ L\"uscher}
{Chiral symmetry and the Yang--Mills gradient flow}
{JHEP \textbf{04} (2013) 123}
{arXiv:1302.5246 | DOI: 10.1007/JHEP04(2013)123}
\whyselected{Third hub of the flow corpus ($\sim$12/99): the fermion flow
defined here is the basis of ringed fermions, smeared bilinears, the
smeared-quasi-PDF program and the flow-dressed CP-violating operators all
present in the library. Original article of the library.}
\physics{The flow is extended to quark fields by the covariant heat
equation
\begin{equation}
  \partial_t\chi = \Delta\,\chi, \qquad \Delta = D_\mu D_\mu,
  \qquad \chi(0)=\psi ,
\end{equation}
whose solution admits a representation as a five-dimensional local theory
(the fermion propagates in one extra dimension). Consequences proved or
constructed here: flowed bilinears such as $S(t)=\bar\chi\chi$ have finite
continuum limits; the small-flow-time limit of $S(t)$ defines the
renormalized scalar density, so the chiral condensate follows from
\begin{equation}
  \lim_{t\to0} Z_S^{-1}(t)\,\langle S(t)\rangle = -\langle\bar\psi\psi\rangle_R ,
\end{equation}
with no additive self-renormalization ambiguity; and O($a$)-improved
lattice actions can be constructed by adding the flow-induced dimension-5
term --- the germ of ``ringed'' fermions later used to keep wave-function
renormalization multiplicative in smeared quasi-PDFs.}
\impact{Monahan--Orginos's smeared quasi-PDFs (Chapter~6) and Shindler's
gradient-flow moments (in the library) take the flowed fermion as their
starting point; the NNLO bilinear-matching papers compute exactly the
coefficients of this expansion.}
\cardend

% ----------------------------------------------------------------------
\section{The Energy--Momentum Tensor from the Flow}
\papercard{H.\ Suzuki}
{Energy--momentum tensor from the Yang--Mills gradient flow}
{Prog.\ Theor.\ Exp.\ Phys.\ \textbf{2013}, 083B03}
{arXiv:1304.0533 | DOI: 10.1093/ptep/ptt059}
\whyselected{Hub of the flow corpus ($\sim$8/99), together with its
companion Makino--Suzuki, PTEP \textbf{2014}, 063B02 [arXiv:1403.4772]
which extends the formula to full QCD with fermions ($\sim$9/99). Both
are original articles of the library.}
\physics{Because flowed products are finite (\ref{eq:flow} above and the
L\"uscher--Weisz theorem), any properly normalized local product of flowed
fields has an unambiguous small-flow-time expansion. Suzuki constructs
gauge-invariant combinations $U_{\mu\nu}(t,x)$ of flowed fields whose
expansion begins with the traceless stress tensor, and derives the master
formula for pure YM:
\begin{equation}
  T_{\mu\nu}(x)
   = \lim_{t\to0}\Big[\mathcal{C}_1(t)\,U_{\mu\nu}^{(1)}(t,x)
      + \mathcal{C}_2(t)\,\delta_{\mu\nu}\,U^{(2)}(t,x)\Big],
\end{equation}
where the coefficients are fixed by matching to dimensional
regularization through the trace identity
$T^\mu{}_\mu=(\beta(g)/2g)\,F^2$. The result bridges lattice and continuum
schemes and turns thermodynamic quantities (trace anomaly, entropy
density) into flow measurements.}
\impact{The Makino--Suzuki extension supplies the fermion terms used by
the library's thermodynamics-oriented flow papers; Harlander--Kluth--Lange
(in the library) pushed the Wilson coefficients to NNLO, making the method
quantitatively competitive. Conceptually, this paper is the template for
\emph{any} flowed-composite observable: define it at finite $t$, expand,
match. Chapter~6 will apply the same template to parton physics.}
\cardend
